\section{System Architecture}
\label{sec:architecture}

The CrowdFund platform follows a hybrid architecture that strategically partitions data between on-chain and off-chain storage to optimise for gas efficiency, data integrity, and user experience.

\subsection{Architectural Overview}

The system comprises three primary layers:

\begin{enumerate}[leftmargin=1.5em]
    \item \textbf{Smart Contract Layer (On-Chain):} The \texttt{CrowdFunding.sol} contract deployed on Ethereum manages all financial logic---campaign lifecycle, contributions, withdrawals, and refunds. All monetary state is stored immutably on-chain.
    \item \textbf{Metadata Layer (Off-Chain):} Campaign descriptive data (titles, descriptions, images) is stored off-chain via IPFS (Pinata) or a local server fallback. The on-chain contract stores only a \texttt{bytes32} IPFS CID as a pointer to this data.
    \item \textbf{Frontend Layer:} A React-based single-page application communicates with the smart contract via ethers.js and MetaMask, providing users with an intuitive interface for all platform interactions.
\end{enumerate}

\subsection{On-Chain Data Model}

All data critical to the financial integrity of the platform is stored directly in the smart contract:

\begin{table}[H]
\centering
\caption{On-Chain Data Fields}
\label{tab:onchain}
\begin{tabularx}{\textwidth}{l l X}
\toprule
\textbf{Field} & \textbf{Type} & \textbf{Purpose} \\
\midrule
Campaign ID        & \texttt{uint256}    & Auto-incrementing unique identifier \\
Creator address    & \texttt{address}    & Wallet address of the campaign creator \\
Goal               & \texttt{uint256}    & Funding target in Wei \\
Deadline           & \texttt{uint256}    & Unix timestamp after which no contributions are accepted \\
Amount raised      & \texttt{uint256}    & Cumulative contributions received in Wei \\
Status             & \texttt{enum}       & Lifecycle state: Active, Successful, Failed, FundsWithdrawn \\
Contributions      & \texttt{mapping}    & Per-contributor amounts: \texttt{campaignId $\Rightarrow$ address $\Rightarrow$ uint256} \\
IPFS CID           & \texttt{bytes32}    & Hash pointer to off-chain metadata \\
\bottomrule
\end{tabularx}
\end{table}

\subsection{Off-Chain Data Strategy}

Storing large textual or media content on-chain is prohibitively expensive. A single \texttt{SSTORE} operation for a 32-byte value costs approximately 20,000 gas, while storing a 256-character string would cost upwards of 160,000 gas. The platform therefore stores the following data off-chain:

\begin{itemize}[leftmargin=1.5em]
    \item Campaign title and full description
    \item Pitch images and videos
    \item Backer reward tiers and fulfilment details
    \item Creator identity and verification documents
    \item Campaign updates and backer communications
\end{itemize}

The off-chain storage pipeline operates as follows:
\begin{enumerate}[leftmargin=1.5em]
    \item The frontend serialises campaign metadata as JSON.
    \item The metadata is uploaded to Pinata (IPFS pinning service) via the backend API.
    \item Pinata returns an IPFS CID (content hash).
    \item The CID is converted to \texttt{bytes32} and stored on-chain alongside the campaign.
    \item When retrieving campaign details, the frontend reads the \texttt{bytes32} CID from the contract and fetches the full metadata from IPFS gateways or the local server fallback.
\end{enumerate}

If Pinata is unavailable (e.g., due to network restrictions), the backend falls back to local filesystem storage, saving JSON files in a \texttt{metadata/} directory and returning the sanitised filename as the CID.

\subsection{Technology Stack}

\begin{table}[H]
\centering
\caption{Technology Stack}
\label{tab:techstack}
\begin{tabularx}{\textwidth}{l l X}
\toprule
\textbf{Layer} & \textbf{Technology} & \textbf{Role} \\
\midrule
Smart Contract  & Solidity \textasciicircum 0.8.24      & Core business logic \\
Framework       & Foundry (Forge, Cast, Anvil)           & Compilation, testing, deployment \\
Security        & OpenZeppelin ReentrancyGuard            & Reentrancy attack prevention \\
Frontend        & React + Vite                            & Single-page application \\
Web3 Bridge     & ethers.js v6                            & Ethereum RPC \& wallet interaction \\
Wallet          & MetaMask                                & User authentication \& transaction signing \\
Off-Chain Store & Pinata (IPFS) / Local Server            & Campaign metadata persistence \\
Backend         & Express.js (Node.js)                    & API proxy for IPFS uploads \\
Deployment      & Vercel                                  & Production hosting (frontend + serverless API) \\
Local Node      & Anvil                                   & Local Ethereum development chain \\
\bottomrule
\end{tabularx}
\end{table}

\subsection{Project Directory Structure}

\begin{lstlisting}[language={}, caption={Project Structure}, label={lst:structure}]
Crowd-Funding/
|-- src/
|   +-- CrowdFunding.sol          # Main smart contract
|-- test/
|   +-- CrowdFunding.t.sol        # Comprehensive test suite
|-- script/
|   +-- deploy.s.sol              # Foundry deployment script
|-- frontend/
|   |-- src/
|   |   |-- components/           # React UI components
|   |   |-- config/contract.js    # ABI + deployed address
|   |   |-- context/Web3Context.jsx
|   |   |-- hooks/useContract.js  # Contract interaction hook
|   |   +-- App.jsx               # Main application
|   |-- api/                      # Vercel serverless functions
|   |-- server.js                 # Express backend (local dev)
|   +-- vercel.json               # Vercel deployment config
|-- lib/
|   |-- forge-std/                # Foundry standard library
|   +-- openzeppelin-contracts/   # OpenZeppelin contracts
|-- foundry.toml                  # Foundry configuration
|-- gas-report.md                 # Gas optimisation report
+-- README.md
\end{lstlisting}
